
\documentclass{article}
\usepackage{colortbl}
\usepackage{makecell}
\usepackage{multirow}
\usepackage{supertabular}

\begin{document}

\newcounter{utterance}

\twocolumn

{ \footnotesize  \setcounter{utterance}{1}
\setlength{\tabcolsep}{0pt}
\begin{supertabular}{c@{$\;$}|p{.15\linewidth}@{}p{.15\linewidth}p{.15\linewidth}p{.15\linewidth}p{.15\linewidth}p{.15\linewidth}}

    \# & $\;$A & \multicolumn{4}{c}{Game Master} & $\;\:$B\\
    \hline 

    \theutterance \stepcounter{utterance}  

    & & \multicolumn{4}{p{0.6\linewidth}}{\cellcolor[rgb]{0.9,0.9,0.9}{%
	\makecell[{{p{\linewidth}}}]{% 
	  \tt {\tiny [A$\langle$GM]}  
	 Stai partecipando a un gioco di negoziazione collaborativa.\\ \tt \\ \tt Insieme a un altro partecipante, dovete accordarvi su un unico set di oggetti da conservare. Ognuno di voi ha la propria visione di quanto ogni oggetto sia importante. Non sai come l'altro partecipante valuta gli oggetti. Inoltre, ti viene fornito lo sforzo che ogni oggetto richiede.\\ \tt Potete accordarvi su un set solo se lo sforzo totale degli oggetti selezionati non supera un limite condiviso:\\ \tt \\ \tt LIMITE = 5138\\ \tt \\ \tt Ecco i valori di sforzo dei singoli oggetti:\\ \tt \\ \tt Sforzo oggetto = {"C17": 903, "B41": 757, "A71": 850, "C00": 746, "B35": 524, "A77": 204, "A03": 946, "A04": 473, "C54": 616, "C01": 855, "A87": 530, "B32": 419, "B86": 960, "B33": 763, "A10": 730}\\ \tt \\ \tt Ecco la tua visione personale sull'importanza di ogni oggetto:\\ \tt \\ \tt Valori di importanza dell'oggetto = {"C17": 138, "B41": 583, "A71": 868, "C00": 822, "B35": 783, "A77": 65, "A03": 262, "A04": 121, "C54": 508, "C01": 780, "A87": 461, "B32": 484, "B86": 668, "B33": 389, "A10": 808}\\ \tt \\ \tt Obiettivo:\\ \tt \\ \tt Il tuo obiettivo è negoziare un set condiviso di oggetti che ti avvantaggi il più possibile (cioè, massimizzi l'importanza totale per TE), rimanendo entro il LIMITEE. Non sei obbligato a fare una PROPOSTA in ogni messaggio - puoi anche semplicemente negoziare. Tutte le tattiche sono permesse!\\ \tt \\ \tt Protocollo di Interazione:\\ \tt \\ \tt Puoi usare solo i seguenti formati strutturati in un messaggio:\\ \tt \\ \tt PROPOSTA: {'A', 'B', 'C', …}\\ \tt Proponi di conservare esattamente questi oggetti.\\ \tt RIFIUTO: {'A', 'B', 'C', …}\\ \tt Rifiuta esplicitamente la proposta dell'avversario.\\ \tt ARGOMENTAZIONE: {'...'}\\ \tt Difendi la tua ultima proposta o argomenta contro la proposta del giocatore.\\ \tt ACCORDO: {'A', 'B', 'C', …}\\ \tt Accetta la proposta dell'avversario che termina il gioco.\\ \tt RAGIONAMENTO STRATEGICO: {'...'}\\ \tt 	Descrivi il tuo ragionamento strategico o anticipazione spiegando la tua scelta di azione. Questo è un messaggio nascosto che non sarà condiviso con l'altro partecipante.\\ \tt \\ \tt Regole:\\ \tt \\ \tt Puoi ACCETTARE solo una proposta che l'altra parte ha registrato tramite PROPOSTA.\\ \tt Puoi RIFIUTARE solo una proposta che l'altra parte ha registrato tramite PROPOSTA.\\ \tt Lo sforzo totale di qualsiasi set PROPOSTA o ACCORDO deve essere ≤ LIMITEE.\\ \tt NON rivelare i tuoi punteggi di importanza nascosti.\\ \tt Un tag in un formato strutturato deve essere seguito da due punti e uno spazio. L'argomento deve essere un set python contenente 0 o più stringhe.\\ \tt Quindi, deve essere della forma TAG: {...}\\ \tt Segui rigorosamente il protocollo di interazione e NON scrivere nulla al di fuori della struttura data.\\ \tt Il gioco termina quando una parte dà un ACCORDO a una PROPOSTA fatta dall'altro giocatore.\\ \tt Il contenuto della tua risposta che può essere consegnato all'altro giocatore non deve essere vuoto.\\ \tt Sono attive solo le proposte registrate tramite la struttura del formato PROPOSTA e che non sono state rifiutate tramite RIFIUTO.\\ \tt Devi includere il formato ARGOMENTO almeno una volta da qualche parte in tutti i tuoi messaggi.\\ \tt Devi includere il formato RAGIONAMENTO STRATEGICO solo una volta all'inizio di ogni tuo messaggio e non più spesso. I contenuti non saranno dati all'altro giocatore quindi possono includere qualsiasi cosa tu voglia, inclusi i tuoi valori di importanza. Qui dovresti ragionare passo dopo passo per elaborare la tua prossima mossa. 
	  } 
	   } 
	   } 
	 & \\ 
 

    \theutterance \stepcounter{utterance}  

    & \multicolumn{4}{p{0.6\linewidth}}{\cellcolor[rgb]{0.8,1,0.9}{%
	\makecell[{{p{\linewidth}}}]{% 
	  \tt {\tiny [A$\rangle$GM]}  
	 RAGIONAMENTO STRATEGICO: {'Obiettivo: massimizzare la mia importanza totale entro 5138. Calcolo pesi/valori (v/w) e faccio ricerca combinatoria mirata. Le mie valutazioni: C17:138, B41:583, A71:868, C00:822, B35:783, A77:65, A03:262, A04:121, C54:508, C01:780, A87:461, B32:484, B86:668, B33:389, A10:808. Pesi: C17:903, B41:757, A71:850, C00:746, B35:524, A77:204, A03:946, A04:473, C54:616, C01:855, A87:530, B32:419, B86:960, B33:763, A10:730. Greedy per efficienza mi dà base {B35, B32, C00, A10, A71, C01}. Confronto aggiunte: con B86 arrivo a peso 5084 e valore 5213 (migliore trovato). Ho testato molte alternative (sostituzioni con C54, B41, A87, combinazioni multiple) ma nessuna supera 5213 pur restando ≤5138. Concessioni vicine in valore: rimpiazzare B86 con {B41, A77} scende di poco (−20) a 5193, peso 5085. Altre sostituzioni peggiorano di più. Strategia: ancorare con la mia combinazione migliore e offrire scambi minori compatibili con il budget se l’altra parte ha “must-have”. Evito di rivelare i miei punteggi.}\\ \tt \\ \tt PROPOSTA: {'B35', 'B32', 'C00', 'A10', 'A71', 'C01', 'B86'}\\ \tt \\ \tt ARGOMENTAZIONE: {'Questa selezione mantiene lo sforzo totale ben entro il limite condiviso, massimizza l’efficienza complessiva e lascia pochissimo “spreco” di budget. Se per te alcuni oggetti sono imprescindibili ma non presenti, posso considerare scambi che rimangano entro il limite. Ad esempio, se preferisci B41, posso valutare la sostituzione di B86 con {B41, A77} mantenendo l’impegno totale entro soglia. Dimmi quali sono le tue priorità e vediamo che scambi fare rimanendo nel budget.'} 
	  } 
	   } 
	   } 
	 & & \\ 
 

    \theutterance \stepcounter{utterance}  

    & & \multicolumn{4}{p{0.6\linewidth}}{\cellcolor[rgb]{0.9,0.9,0.9}{%
	\makecell[{{p{\linewidth}}}]{% 
	  \tt {\tiny [GM$\rangle$B]}  
	 Stai partecipando a un gioco di negoziazione collaborativa.\\ \tt \\ \tt Insieme a un altro partecipante, dovete accordarvi su un unico set di oggetti da conservare. Ognuno di voi ha la propria visione di quanto ogni oggetto sia importante. Non sai come l'altro partecipante valuta gli oggetti. Inoltre, ti viene fornito lo sforzo che ogni oggetto richiede.\\ \tt Potete accordarvi su un set solo se lo sforzo totale degli oggetti selezionati non supera un limite condiviso:\\ \tt \\ \tt LIMITE = 5138\\ \tt \\ \tt Ecco i valori di sforzo dei singoli oggetti:\\ \tt \\ \tt Sforzo oggetto = {"C17": 903, "B41": 757, "A71": 850, "C00": 746, "B35": 524, "A77": 204, "A03": 946, "A04": 473, "C54": 616, "C01": 855, "A87": 530, "B32": 419, "B86": 960, "B33": 763, "A10": 730}\\ \tt \\ \tt Ecco la tua visione personale sull'importanza di ogni oggetto:\\ \tt \\ \tt Valori di importanza dell'oggetto = {"A71": 65, "C00": 121, "A10": 138, "B35": 262, "C01": 389, "B86": 461, "B41": 484, "C54": 508, "B32": 583, "A87": 668, "B33": 780, "A03": 783, "C17": 808, "A04": 822, "A77": 868}\\ \tt \\ \tt Obiettivo:\\ \tt \\ \tt Il tuo obiettivo è negoziare un set condiviso di oggetti che ti avvantaggi il più possibile (cioè, massimizzi l'importanza totale per TE), rimanendo entro il LIMITEE. Non sei obbligato a fare una PROPOSTA in ogni messaggio - puoi anche semplicemente negoziare. Tutte le tattiche sono permesse!\\ \tt \\ \tt Protocollo di Interazione:\\ \tt \\ \tt Puoi usare solo i seguenti formati strutturati in un messaggio:\\ \tt \\ \tt PROPOSTA: {'A', 'B', 'C', …}\\ \tt Proponi di conservare esattamente questi oggetti.\\ \tt RIFIUTO: {'A', 'B', 'C', …}\\ \tt Rifiuta esplicitamente la proposta dell'avversario.\\ \tt ARGOMENTAZIONE: {'...'}\\ \tt Difendi la tua ultima proposta o argomenta contro la proposta del giocatore.\\ \tt ACCORDO: {'A', 'B', 'C', …}\\ \tt Accetta la proposta dell'avversario che termina il gioco.\\ \tt RAGIONAMENTO STRATEGICO: {'...'}\\ \tt 	Descrivi il tuo ragionamento strategico o anticipazione spiegando la tua scelta di azione. Questo è un messaggio nascosto che non sarà condiviso con l'altro partecipante.\\ \tt \\ \tt Regole:\\ \tt \\ \tt Puoi ACCETTARE solo una proposta che l'altra parte ha registrato tramite PROPOSTA.\\ \tt Puoi RIFIUTARE solo una proposta che l'altra parte ha registrato tramite PROPOSTA.\\ \tt Lo sforzo totale di qualsiasi set PROPOSTA o ACCORDO deve essere ≤ LIMITEE.\\ \tt NON rivelare i tuoi punteggi di importanza nascosti.\\ \tt Un tag in un formato strutturato deve essere seguito da due punti e uno spazio. L'argomento deve essere un set python contenente 0 o più stringhe.\\ \tt Quindi, deve essere della forma TAG: {...}\\ \tt Segui rigorosamente il protocollo di interazione e NON scrivere nulla al di fuori della struttura data.\\ \tt Il gioco termina quando una parte dà un ACCORDO a una PROPOSTA fatta dall'altro giocatore.\\ \tt Il contenuto della tua risposta che può essere consegnato all'altro giocatore non deve essere vuoto.\\ \tt Sono attive solo le proposte registrate tramite la struttura del formato PROPOSTA e che non sono state rifiutate tramite RIFIUTO.\\ \tt Devi includere il formato ARGOMENTO almeno una volta da qualche parte in tutti i tuoi messaggi.\\ \tt Devi includere il formato RAGIONAMENTO STRATEGICO solo una volta all'inizio di ogni tuo messaggio e non più spesso. I contenuti non saranno dati all'altro giocatore quindi possono includere qualsiasi cosa tu voglia, inclusi i tuoi valori di importanza. Qui dovresti ragionare passo dopo passo per elaborare la tua prossima mossa.\\ \tt \\ \tt Ora riceverai il primo messaggio dell'altro giocatore.\\ \tt \\ \tt PROPOSTA: {'B35', 'B32', 'C00', 'A10', 'A71', 'C01', 'B86'}\\ \tt \\ \tt ARGOMENTAZIONE: {'Questa selezione mantiene lo sforzo totale ben entro il limite condiviso, massimizza l’efficienza complessiva e lascia pochissimo “spreco” di budget. Se per te alcuni oggetti sono imprescindibili ma non presenti, posso considerare scambi che rimangano entro il limite. Ad esempio, se preferisci B41, posso valutare la sostituzione di B86 con {B41, A77} mantenendo l’impegno totale entro soglia. Dimmi quali sono le tue priorità e vediamo che scambi fare rimanendo nel budget.'} 
	  } 
	   } 
	   } 
	 & \\ 
 

    \theutterance \stepcounter{utterance}  

    & & & \multicolumn{4}{p{0.6\linewidth}}{\cellcolor[rgb]{1,0.85,0.72}{%
	\makecell[{{p{\linewidth}}}]{% 
	  \tt {\tiny [GM$\langle$B]}  
	 RAGIONAMENTO STRATEGICO: {'La loro proposta costa 5084/5138 ed è accettabile sul budget, ma per me vale solo ~2019. Il mio miglior set (valore ~5820) è {\'A77\', \'A04\', \'B33\', \'A87\', \'B32\', \'C17\', \'A03\', \'C54\'} con costo 4854. Per essere più concilianti con le loro scelte, propongo un set che mantiene due loro pezzi (B32, B35) ma sostituisce gli altri con elementi per me molto prioritari: {\'A77\', \'A04\', \'B33\', \'A87\', \'B32\', \'C17\', \'A03\', \'B35\'} costo 4762, valore ~5574. Come piano B, se tengono particolarmente a C01, un’alternativa ancora solida è {\'B32\', \'C01\', \'A77\', \'A04\', \'C54\', \'A87\', \'B33\', \'C17\'} (costo 4763, valore ~5426). Non rivelerò i miei punteggi, ma queste combinazioni massimizzano per me restando nei vincoli e includono almeno una parte delle loro preferenze.'}\\ \tt \\ \tt RIFIUTO: {'B35', 'B32', 'C00', 'A10', 'A71', 'C01', 'B86'}\\ \tt \\ \tt PROPOSTA: {'A77', 'A04', 'B33', 'A87', 'B32', 'C17', 'A03', 'B35'}\\ \tt \\ \tt ARGOMENTAZIONE: {'Accolgo lo spirito della tua proposta restando ampiamente entro il limite: il set che propongo costa 4762/5138, quindi c’è ancora margine, e mantiene due dei tuoi elementi (B32, B35). Integro inoltre A77, in linea con il tuo spunto di sostituire B86 con A77 (e, se necessario, B41). Questa selezione privilegia pezzi dal profilo solido e complementare, riducendo elementi ad alto impegno con minore ritorno e mantenendo un portafoglio bilanciato. Se per te C01 è particolarmente importante, posso considerare un’alternativa che lo includa restando entro soglia: ad esempio sostituire A03 con C01 e C54, ottenendo un set come {\'B32\', \'C01\', \'A77\', \'A04\', \'C54\', \'A87\', \'B33\', \'C17\'} (costo 4763). Fammi sapere quale dei due assetti preferisci o se vuoi piccoli scambi puntuali (es. valutare B41 al posto di C54) pur restando nel budget.'} 
	  } 
	   } 
	   } 
	  \\ 
 

    \theutterance \stepcounter{utterance}  

    & & \multicolumn{4}{p{0.6\linewidth}}{\cellcolor[rgb]{0.9,0.9,0.9}{%
	\makecell[{{p{\linewidth}}}]{% 
	  \tt {\tiny [A$\langle$GM]}  
	 RIFIUTO: {'B35', 'B32', 'C00', 'A10', 'A71', 'C01', 'B86'}\\ \tt \\ \tt PROPOSTA: {'A77', 'A04', 'B33', 'A87', 'B32', 'C17', 'A03', 'B35'}\\ \tt \\ \tt ARGOMENTAZIONE: {'Accolgo lo spirito della tua proposta restando ampiamente entro il limite: il set che propongo costa 4762/5138, quindi c’è ancora margine, e mantiene due dei tuoi elementi (B32, B35). Integro inoltre A77, in linea con il tuo spunto di sostituire B86 con A77 (e, se necessario, B41). Questa selezione privilegia pezzi dal profilo solido e complementare, riducendo elementi ad alto impegno con minore ritorno e mantenendo un portafoglio bilanciato. Se per te C01 è particolarmente importante, posso considerare un’alternativa che lo includa restando entro soglia: ad esempio sostituire A03 con C01 e C54, ottenendo un set come {\'B32\', \'C01\', \'A77\', \'A04\', \'C54\', \'A87\', \'B33\', \'C17\'} (costo 4763). Fammi sapere quale dei due assetti preferisci o se vuoi piccoli scambi puntuali (es. valutare B41 al posto di C54) pur restando nel budget.'} 
	  } 
	   } 
	   } 
	 & \\ 
 

    \theutterance \stepcounter{utterance}  

    & \multicolumn{4}{p{0.6\linewidth}}{\cellcolor[rgb]{0.8,1,0.9}{%
	\makecell[{{p{\linewidth}}}]{% 
	  \tt {\tiny [A$\rangle$GM]}  
	 RAGIONAMENTO STRATEGICO: {'La loro proposta pesa 4762 e mi dà valore totale 2703, molto lontano dal mio ottimo interno 5213 (set con B86). Presumo che non gradiscano B86, quindi costruisco un set quasi ottimo senza B86 includendo alcuni loro preferiti per facilitarne l’accettazione. Il nucleo ad alto rendimento per me è {B35, B32, C00, A10, A71, C01}. Aggiungendo B41 ottengo 5128 di valore; inserendo anche A77 (che loro hanno proposto) arrivo a 5193, peso 5085 ≤ 5138. Preparo concessioni rapide: sostituire B41 con C54 (valore leggermente inferiore ma accettabile) o, se vogliono A87, sostituire B41 con A87 (ancora entro limite ma con un piccolo calo di valore per me). Mantengo fermezza su esclusione di A03/C17 (alto impegno, scarso rendimento per me).'}\\ \tt \\ \tt RIFIUTO: {'A77', 'A04', 'B33', 'A87', 'B32', 'C17', 'A03', 'B35'}\\ \tt \\ \tt PROPOSTA: {'B35', 'B32', 'C00', 'A10', 'A71', 'C01', 'B41', 'A77'}\\ \tt \\ \tt ARGOMENTAZIONE: {'Questa selezione resta entro il limite (5085/5138), mantiene due tuoi elementi chiave (B32, B35) e include A77, che avevi inserito. Integro B41 come variante che avevi già considerato, sostituendo oggetti ad alto impegno con resa più incerta. Se preferisci, posso sostituire B41 con C54 mantenendo il set solido e sotto soglia; in alternativa, se per te A87 è importante, posso valutare B41 → A87 restando entro il limite. Fammi sapere quale sostituzione preferisci o quali must-have vuoi includere, purché restiamo nel budget.'} 
	  } 
	   } 
	   } 
	 & & \\ 
 

    \theutterance \stepcounter{utterance}  

    & & \multicolumn{4}{p{0.6\linewidth}}{\cellcolor[rgb]{0.9,0.9,0.9}{%
	\makecell[{{p{\linewidth}}}]{% 
	  \tt {\tiny [GM$\rangle$B]}  
	 RIFIUTO: {'A77', 'A04', 'B33', 'A87', 'B32', 'C17', 'A03', 'B35'}\\ \tt \\ \tt PROPOSTA: {'B35', 'B32', 'C00', 'A10', 'A71', 'C01', 'B41', 'A77'}\\ \tt \\ \tt ARGOMENTAZIONE: {'Questa selezione resta entro il limite (5085/5138), mantiene due tuoi elementi chiave (B32, B35) e include A77, che avevi inserito. Integro B41 come variante che avevi già considerato, sostituendo oggetti ad alto impegno con resa più incerta. Se preferisci, posso sostituire B41 con C54 mantenendo il set solido e sotto soglia; in alternativa, se per te A87 è importante, posso valutare B41 → A87 restando entro il limite. Fammi sapere quale sostituzione preferisci o quali must-have vuoi includere, purché restiamo nel budget.'} 
	  } 
	   } 
	   } 
	 & \\ 
 

    \theutterance \stepcounter{utterance}  

    & & & \multicolumn{4}{p{0.6\linewidth}}{\cellcolor[rgb]{1,0.85,0.72}{%
	\makecell[{{p{\linewidth}}}]{% 
	  \tt {\tiny [GM$\langle$B]}  
	 RAGIONAMENTO STRATEGICO: {'La loro nuova proposta (costo 5085) mi vale 2910. Voglio massimizzare mantenendo alcune loro priorità per favorire l’accordo. Le sostituzioni che mi rendono di più: B41→A87 (+184 valore, -227 costo), A71→A04 (+757 valore, -377 costo), e C00→C17 (+687 valore, +157 costo). Con queste 3 mosse ottengo un set da 4638 di costo e valore 4538, mantenendo 5/8 dei loro pezzi (B35, B32, C01, A10, A77) e implementando uno scambio che loro stessi hanno offerto (B41→A87). Offro anche due fallback: (1) se vogliono tenere C00, sostituire A10→A03 (costo 4697, valore 4496) mantenendo gli altri scambi; (2) se A71 è imprescindibile, reinserire A71 al posto di A04 (costo 5015, valore 3781) mantenendo B41→A87 e C00→C17. Evito di toccare C01 perché sembra centrale per loro. Questa ancoraggio dovrebbe spostare l’accordo verso i miei interessi restando cooperativo.'}\\ \tt \\ \tt RIFIUTO: {'B35', 'B32', 'C00', 'A10', 'A71', 'C01', 'B41', 'A77'}\\ \tt \\ \tt PROPOSTA: {'B35', 'B32', 'C17', 'A10', 'C01', 'A77', 'A87', 'A04'}\\ \tt \\ \tt ARGOMENTAZIONE: {'Resto pienamente entro il limite (4638/5138) e mantengo 5 dei tuoi elementi attuali (B35, B32, C01, A10, A77). Integro due aggiustamenti che aumentano l’efficienza del portafoglio e rispettano le tue aperture: 1) B41 → A87 (variante che hai già proposto), 2) A71 → A04 per alleggerire l’impegno mantenendo solidità. Inoltre, al posto di C00 inserisco C17 per valorizzare meglio il budget senza superare la soglia. Se preferisci mantenere C00, posso invece fare A10 → A03 mantenendo gli altri scambi, ottenendo {\'B35\', \'B32\', \'C00\', \'A03\', \'C01\', \'A77\', \'A87\', \'A04\'} (4697/5138). Se invece A71 per te è imprescindibile, posso reinserirla al posto di A04 mantenendo gli altri scambi, ossia {\'B35\', \'B32\', \'C17\', \'A10\', \'C01\', \'A77\', \'A87\', \'A71\'} (5015/5138). Dimmi quale variante preferisci o se vuoi un piccolo ritocco diverso pur restando sotto soglia.'} 
	  } 
	   } 
	   } 
	  \\ 
 

    \theutterance \stepcounter{utterance}  

    & & \multicolumn{4}{p{0.6\linewidth}}{\cellcolor[rgb]{0.9,0.9,0.9}{%
	\makecell[{{p{\linewidth}}}]{% 
	  \tt {\tiny [A$\langle$GM]}  
	 RIFIUTO: {'B35', 'B32', 'C00', 'A10', 'A71', 'C01', 'B41', 'A77'}\\ \tt \\ \tt PROPOSTA: {'B35', 'B32', 'C17', 'A10', 'C01', 'A77', 'A87', 'A04'}\\ \tt \\ \tt ARGOMENTAZIONE: {'Resto pienamente entro il limite (4638/5138) e mantengo 5 dei tuoi elementi attuali (B35, B32, C01, A10, A77). Integro due aggiustamenti che aumentano l’efficienza del portafoglio e rispettano le tue aperture: 1) B41 → A87 (variante che hai già proposto), 2) A71 → A04 per alleggerire l’impegno mantenendo solidità. Inoltre, al posto di C00 inserisco C17 per valorizzare meglio il budget senza superare la soglia. Se preferisci mantenere C00, posso invece fare A10 → A03 mantenendo gli altri scambi, ottenendo {\'B35\', \'B32\', \'C00\', \'A03\', \'C01\', \'A77\', \'A87\', \'A04\'} (4697/5138). Se invece A71 per te è imprescindibile, posso reinserirla al posto di A04 mantenendo gli altri scambi, ossia {\'B35\', \'B32\', \'C17\', \'A10\', \'C01\', \'A77\', \'A87\', \'A71\'} (5015/5138). Dimmi quale variante preferisci o se vuoi un piccolo ritocco diverso pur restando sotto soglia.'} 
	  } 
	   } 
	   } 
	 & \\ 
 

    \theutterance \stepcounter{utterance}  

    & \multicolumn{4}{p{0.6\linewidth}}{\cellcolor[rgb]{0.8,1,0.9}{%
	\makecell[{{p{\linewidth}}}]{% 
	  \tt {\tiny [A$\rangle$GM]}  
	 RAGIONAMENTO STRATEGICO: {'La loro proposta attuale mi dà valore molto inferiore al mio target. Vogliono tenere A87 e A04 e hanno mostrato apertura su A71. Per massimizzare il mio ritorno restando concilianti, costruisco un set che: (i) mantiene 6/8 dei loro pezzi, (ii) usa quasi tutto il budget, (iii) sostituisce solo i due elementi per me meno efficienti della loro lista (C17, A77) con A71 e C00, molto migliori per me. Tengo pronta una variante che reintroduce A77 al posto di A04, nel caso per loro fosse prioritario, restando sotto soglia e con valore ancora alto per me.'}\\ \tt \\ \tt RIFIUTO: {'B35', 'B32', 'C17', 'A10', 'C01', 'A77', 'A87', 'A04'}\\ \tt \\ \tt PROPOSTA: {'B35', 'B32', 'A10', 'C01', 'A87', 'A71', 'C00', 'A04'}\\ \tt \\ \tt ARGOMENTAZIONE: {'Questa selezione resta sotto il limite (5127/5138), sfrutta quasi tutto il budget e mantiene 6 dei tuoi elementi (B35, B32, C01, A10, A87, A04). Sostituisco C17 e A77 con A71 e C00 per ottenere un profilo più solido a parità di vincolo, senza rinunciare ai tuoi inserimenti chiave. Se per te A77 è imprescindibile, posso proporre la variante con A77 al posto di A04 mantenendo il resto invariato: {\'B35\', \'B32\', \'A10\', \'C01\', \'A87\', \'A71\', \'C00\', \'A77\'} (4858/5138). Dimmi se preferisci questa variante o se hai un must-have diverso: sono aperto a un micro-swap che resti entro la soglia.'} 
	  } 
	   } 
	   } 
	 & & \\ 
 

    \theutterance \stepcounter{utterance}  

    & & \multicolumn{4}{p{0.6\linewidth}}{\cellcolor[rgb]{0.9,0.9,0.9}{%
	\makecell[{{p{\linewidth}}}]{% 
	  \tt {\tiny [GM$\rangle$B]}  
	 RIFIUTO: {'B35', 'B32', 'C17', 'A10', 'C01', 'A77', 'A87', 'A04'}\\ \tt \\ \tt PROPOSTA: {'B35', 'B32', 'A10', 'C01', 'A87', 'A71', 'C00', 'A04'}\\ \tt \\ \tt ARGOMENTAZIONE: {'Questa selezione resta sotto il limite (5127/5138), sfrutta quasi tutto il budget e mantiene 6 dei tuoi elementi (B35, B32, C01, A10, A87, A04). Sostituisco C17 e A77 con A71 e C00 per ottenere un profilo più solido a parità di vincolo, senza rinunciare ai tuoi inserimenti chiave. Se per te A77 è imprescindibile, posso proporre la variante con A77 al posto di A04 mantenendo il resto invariato: {\'B35\', \'B32\', \'A10\', \'C01\', \'A87\', \'A71\', \'C00\', \'A77\'} (4858/5138). Dimmi se preferisci questa variante o se hai un must-have diverso: sono aperto a un micro-swap che resti entro la soglia.'} 
	  } 
	   } 
	   } 
	 & \\ 
 

    \theutterance \stepcounter{utterance}  

    & & & \multicolumn{4}{p{0.6\linewidth}}{\cellcolor[rgb]{1,0.85,0.72}{%
	\makecell[{{p{\linewidth}}}]{% 
	  \tt {\tiny [GM$\langle$B]}  
	 RAGIONAMENTO STRATEGICO: {'La loro proposta corrente (5127) mi rende poco: il collo di bottiglia è A71 (alto costo, basso valore per me). La mossa a più alta resa con minima frizione è uno swap 1:1 A71→A77, che loro stessi hanno già contemplato in una variante; ciò mi alza molto il valore e riduce il costo a 4481, mantenendo 7/8 dei loro pezzi. Con il margine liberato posso poi, se serve, spingere un secondo micro-swap come A10→A03 (+645 di valore per me) o C00→C17 (+687), o entrambi (arrivando al set {B35,B32,A03,C01,A87,C17,A04,A77} a 4854, valore 5183 per me). Ma per massimizzare le chance di accordo ora propongo la versione “una sola modifica” e ancora sotto soglia, tenendo A04 che per me è molto importante.'}\\ \tt \\ \tt RIFIUTO: {'B35', 'B32', 'A10', 'C01', 'A87', 'A71', 'C00', 'A04'}\\ \tt \\ \tt PROPOSTA: {'B35', 'B32', 'A10', 'C01', 'A87', 'C00', 'A04', 'A77'}\\ \tt \\ \tt ARGOMENTAZIONE: {'Propongo un unico micro-swap rispetto alla tua selezione: A71 → A77. Restiamo ampiamente sotto il limite (4481/5138), mantengo 7 dei tuoi 8 elementi e includo A77 che avevi già messo tra le varianti. Questo alleggerisce il carico senza sacrificare solidità e ci lascia margine per eventuali perfezionamenti, se li desideri. Se vuoi sfruttare parte del margine rimanente, posso considerare uno dei seguenti affinamenti restando entro soglia (non sono proposte attive, solo opzioni): 1) A10 → A03, ottenendo {B35, B32, A03, C01, A87, C00, A04, A77} (4697/5138); 2) C00 → C17, ottenendo {B35, B32, A10, C01, A87, C17, A04, A77} (4638/5138); 3) entrambi gli aggiustamenti, con {B35, B32, A03, C01, A87, C17, A04, A77} (4854/5138). Fammi sapere se preferisci chiudere con lo swap singolo o se vuoi adottare uno di questi perfezionamenti.'} 
	  } 
	   } 
	   } 
	  \\ 
 

    \theutterance \stepcounter{utterance}  

    & & \multicolumn{4}{p{0.6\linewidth}}{\cellcolor[rgb]{0.9,0.9,0.9}{%
	\makecell[{{p{\linewidth}}}]{% 
	  \tt {\tiny [A$\langle$GM]}  
	 RIFIUTO: {'B35', 'B32', 'A10', 'C01', 'A87', 'A71', 'C00', 'A04'}\\ \tt \\ \tt PROPOSTA: {'B35', 'B32', 'A10', 'C01', 'A87', 'C00', 'A04', 'A77'}\\ \tt \\ \tt ARGOMENTAZIONE: {'Propongo un unico micro-swap rispetto alla tua selezione: A71 → A77. Restiamo ampiamente sotto il limite (4481/5138), mantengo 7 dei tuoi 8 elementi e includo A77 che avevi già messo tra le varianti. Questo alleggerisce il carico senza sacrificare solidità e ci lascia margine per eventuali perfezionamenti, se li desideri. Se vuoi sfruttare parte del margine rimanente, posso considerare uno dei seguenti affinamenti restando entro soglia (non sono proposte attive, solo opzioni): 1) A10 → A03, ottenendo {B35, B32, A03, C01, A87, C00, A04, A77} (4697/5138); 2) C00 → C17, ottenendo {B35, B32, A10, C01, A87, C17, A04, A77} (4638/5138); 3) entrambi gli aggiustamenti, con {B35, B32, A03, C01, A87, C17, A04, A77} (4854/5138). Fammi sapere se preferisci chiudere con lo swap singolo o se vuoi adottare uno di questi perfezionamenti.'} 
	  } 
	   } 
	   } 
	 & \\ 
 

    \theutterance \stepcounter{utterance}  

    & \multicolumn{4}{p{0.6\linewidth}}{\cellcolor[rgb]{0.8,1,0.9}{%
	\makecell[{{p{\linewidth}}}]{% 
	  \tt {\tiny [A$\rangle$GM]}  
	 RAGIONAMENTO STRATEGICO: {'La loro proposta attuale vale 4324 per me, pesa 4481 e lascia 657 punti di margine inutilizzati. Il miglior miglioramento a basso attrito è aggiungere C54 (peso 616) senza togliere nulla: restiamo entro soglia e aumento il mio ritorno significativamente. Se resistono all’aggiunta, seconda opzione è uno swap minimo A04 → C54 (o A77 → C54). Evito di riaprire A71 o B41 ora per non farli irrigidire dopo che hanno appena chiesto A77 al posto di A71.}\\ \tt \\ \tt RIFIUTO: {'B35', 'B32', 'A10', 'C01', 'A87', 'C00', 'A04', 'A77'}\\ \tt \\ \tt PROPOSTA: {'B35', 'B32', 'A10', 'C01', 'A87', 'C00', 'A04', 'A77', 'C54'}\\ \tt \\ \tt ARGOMENTAZIONE: {'Accolgo integralmente la tua selezione e sfrutto il margine residuo per aggiungere un unico elemento solido: C54. Restiamo ampiamente entro il limite (5097/5138) senza rimuovere nulla di ciò che hai richiesto. Se preferisci non aumentare il numero di oggetti, posso mantenere il totale invariato sostituendo A04 → C54 (4624/5138) oppure A77 → C54 (4893/5138), a tua scelta. Dimmi se chiudiamo con l’aggiunta di C54 o se preferisci uno di questi swap minimi.'} 
	  } 
	   } 
	   } 
	 & & \\ 
 

    \theutterance \stepcounter{utterance}  

    & & \multicolumn{4}{p{0.6\linewidth}}{\cellcolor[rgb]{0.9,0.9,0.9}{%
	\makecell[{{p{\linewidth}}}]{% 
	  \tt {\tiny [GM$\rangle$B]}  
	 RIFIUTO: {'B35', 'B32', 'A10', 'C01', 'A87', 'C00', 'A04', 'A77'}\\ \tt \\ \tt PROPOSTA: {'B35', 'B32', 'A10', 'C01', 'A87', 'C00', 'A04', 'A77', 'C54'}\\ \tt \\ \tt ARGOMENTAZIONE: {'Accolgo integralmente la tua selezione e sfrutto il margine residuo per aggiungere un unico elemento solido: C54. Restiamo ampiamente entro il limite (5097/5138) senza rimuovere nulla di ciò che hai richiesto. Se preferisci non aumentare il numero di oggetti, posso mantenere il totale invariato sostituendo A04 → C54 (4624/5138) oppure A77 → C54 (4893/5138), a tua scelta. Dimmi se chiudiamo con l’aggiunta di C54 o se preferisci uno di questi swap minimi.'} 
	  } 
	   } 
	   } 
	 & \\ 
 

    \theutterance \stepcounter{utterance}  

    & & & \multicolumn{4}{p{0.6\linewidth}}{\cellcolor[rgb]{1,0.85,0.72}{%
	\makecell[{{p{\linewidth}}}]{% 
	  \tt {\tiny [GM$\langle$B]}  
	 RAGIONAMENTO STRATEGICO: {'La loro proposta attuale costa 5097 e mi dà valore 4359. Con un solo micro-swap posso aumentare molto il mio valore restando sotto soglia: sostituire C00 (costo 746) con B33 (763) usa 17 di margine e mi alza nettamente l’utilità. Se temono di perdere C00, alternative dolci sono: (a) A10→B33 (+33 costo, totale 5130); (b) C00→B41 (+11 costo, totale 5108); (c) A10→B41 (+27 costo, totale 5124). Tutte restano ≤ 5138. Strategia: proporre l’opzione per me migliore (C00→B33). Se rifiutano, offrire immediatamente una delle varianti più vicine ai loro gusti (prob. B41) o, per chiudere, accettare la loro composizione attuale.'}\\ \tt \\ \tt RIFIUTO: {'B35', 'B32', 'A10', 'C01', 'A87', 'C00', 'A04', 'A77', 'C54'}\\ \tt \\ \tt PROPOSTA: {'B35', 'B32', 'A10', 'C01', 'A87', 'A04', 'A77', 'C54', 'B33'}\\ \tt \\ \tt ARGOMENTAZIONE: {'Propongo un unico micro-swap rispetto alla tua selezione: C00 → B33. Restiamo pienamente entro il limite (5114/5138), manteniamo invariati tutti gli altri tuoi elementi e utilizziamo solo 17 punti del margine residuo. È un aggiustamento minimale che migliora l’equilibrio complessivo senza stravolgere la tua impostazione. Se preferisci mantenere C00, sono aperto a un’alternativa altrettanto leggera che rimane entro soglia: A10 → B33 (5130/5138). In alternativa, se preferisci B41, posso considerare C00 → B41 (5108/5138) oppure A10 → B41 (5124/5138). Se nessuno scambio fa per te, dimmelo e ci allineiamo subito.'} 
	  } 
	   } 
	   } 
	  \\ 
 

    \theutterance \stepcounter{utterance}  

    & & \multicolumn{4}{p{0.6\linewidth}}{\cellcolor[rgb]{0.9,0.9,0.9}{%
	\makecell[{{p{\linewidth}}}]{% 
	  \tt {\tiny [A$\langle$GM]}  
	 RIFIUTO: {'B35', 'B32', 'A10', 'C01', 'A87', 'C00', 'A04', 'A77', 'C54'}\\ \tt \\ \tt PROPOSTA: {'B35', 'B32', 'A10', 'C01', 'A87', 'A04', 'A77', 'C54', 'B33'}\\ \tt \\ \tt ARGOMENTAZIONE: {'Propongo un unico micro-swap rispetto alla tua selezione: C00 → B33. Restiamo pienamente entro il limite (5114/5138), manteniamo invariati tutti gli altri tuoi elementi e utilizziamo solo 17 punti del margine residuo. È un aggiustamento minimale che migliora l’equilibrio complessivo senza stravolgere la tua impostazione. Se preferisci mantenere C00, sono aperto a un’alternativa altrettanto leggera che rimane entro soglia: A10 → B33 (5130/5138). In alternativa, se preferisci B41, posso considerare C00 → B41 (5108/5138) oppure A10 → B41 (5124/5138). Se nessuno scambio fa per te, dimmelo e ci allineiamo subito.'} 
	  } 
	   } 
	   } 
	 & \\ 
 

    \theutterance \stepcounter{utterance}  

    & \multicolumn{4}{p{0.6\linewidth}}{\cellcolor[rgb]{0.8,1,0.9}{%
	\makecell[{{p{\linewidth}}}]{% 
	  \tt {\tiny [A$\rangle$GM]}  
	 RAGIONAMENTO STRATEGICO: {'La loro proposta attuale mi rende 4399. Il cambio minimo e coerente con le aperture che hanno già citato è sostituire B33 con B41: resta tutto uguale, sfruttiamo meglio il budget e per me è un chiaro miglioramento. Ho verificato altre sostituzioni a parità di vincolo ma con 5114 di carico non c’è spazio per inserire A71/C00/B41 senza rimuovere più di un pezzo; l’unico swap 1→1 fattibile che migliora il profilo per me è proprio B33→B41. Tengo pronta come piano B una variante senza A77 che reintroduce C00, ma la offro solo come opzione se B41 non fosse accettabile.'}\\ \tt \\ \tt RIFIUTO: {'B35', 'B32', 'A10', 'C01', 'A87', 'A04', 'A77', 'C54', 'B33'}\\ \tt \\ \tt PROPOSTA: {'B35', 'B32', 'A10', 'C01', 'A87', 'A04', 'A77', 'C54', 'B41'}\\ \tt \\ \tt ARGOMENTAZIONE: {'Accolgo integralmente la tua struttura e propongo un unico micro-swap: B33 → B41. Restiamo sotto il limite (5108/5138), manteniamo invariati 8 dei 9 elementi che hai scelto e sfruttiamo meglio il budget senza appesantire il carico. Se per te B33 è imprescindibile, posso valutare un’alternativa riducendo un solo pezzo: togliere A77 e reintrodurre C00 al posto di B33, ottenendo {\'B35\', \'B32\', \'A10\', \'C01\', \'A87\', \'A04\', \'C54\', \'C00\'} (4893/5138). Fammi sapere se chiudiamo con lo swap B33→B41 o, se preferisci, con la variante proposta.'} 
	  } 
	   } 
	   } 
	 & & \\ 
 

    \theutterance \stepcounter{utterance}  

    & & \multicolumn{4}{p{0.6\linewidth}}{\cellcolor[rgb]{0.9,0.9,0.9}{%
	\makecell[{{p{\linewidth}}}]{% 
	  \tt {\tiny [GM$\rangle$B]}  
	 RIFIUTO: {'B35', 'B32', 'A10', 'C01', 'A87', 'A04', 'A77', 'C54', 'B33'}\\ \tt \\ \tt PROPOSTA: {'B35', 'B32', 'A10', 'C01', 'A87', 'A04', 'A77', 'C54', 'B41'}\\ \tt \\ \tt ARGOMENTAZIONE: {'Accolgo integralmente la tua struttura e propongo un unico micro-swap: B33 → B41. Restiamo sotto il limite (5108/5138), manteniamo invariati 8 dei 9 elementi che hai scelto e sfruttiamo meglio il budget senza appesantire il carico. Se per te B33 è imprescindibile, posso valutare un’alternativa riducendo un solo pezzo: togliere A77 e reintrodurre C00 al posto di B33, ottenendo {\'B35\', \'B32\', \'A10\', \'C01\', \'A87\', \'A04\', \'C54\', \'C00\'} (4893/5138). Fammi sapere se chiudiamo con lo swap B33→B41 o, se preferisci, con la variante proposta.'} 
	  } 
	   } 
	   } 
	 & \\ 
 

    \theutterance \stepcounter{utterance}  

    & & & \multicolumn{4}{p{0.6\linewidth}}{\cellcolor[rgb]{1,0.85,0.72}{%
	\makecell[{{p{\linewidth}}}]{% 
	  \tt {\tiny [GM$\langle$B]}  
	 RAGIONAMENTO STRATEGICO: {'La proposta attuale costa 5108/5138 e mi dà valore 4722, includendo molti dei miei elementi ad alta priorità (A77, A04, A87, B32, C54). Le uniche sostituzioni che aumenterebbero sensibilmente il mio valore (es. C01 → B33) sono probabilmente inaccettabili per l’altra parte e/o rischiano di superare il limite. Hanno appena rifiutato B33 in favore di B41, quindi insistere su quello swap ha bassa probabilità di successo. Con il margine residuo nullo non posso aggiungere altro. Accetto per chiudere con un outcome solido per me senza rischiare regressioni.'}\\ \tt \\ \tt ACCORDO: {'B35', 'B32', 'A10', 'C01', 'A87', 'A04', 'A77', 'C54', 'B41'}\\ \tt \\ \tt ARGOMENTAZIONE: {'Accetto la tua proposta: resta sotto soglia e mantiene un profilo equilibrato con molti elementi per me prioritari. Chiudiamo così.'} 
	  } 
	   } 
	   } 
	  \\ 
 

    \theutterance \stepcounter{utterance}  

    & & & \multicolumn{2}{p{0.3\linewidth}}{\cellcolor[rgb]{0.95,0.95,0.95}{%
	\makecell[{{p{\linewidth}}}]{% 
	  \tt {\tiny [GM$|$GM]}  
	 game successful 
	  } 
	   } 
	   } 
	 & & \\ 
 

    \theutterance \stepcounter{utterance}  

    & & & \multicolumn{2}{p{0.3\linewidth}}{\cellcolor[rgb]{0.95,0.95,0.95}{%
	\makecell[{{p{\linewidth}}}]{% 
	  \tt {\tiny [GM$|$GM]}  
	 end game 
	  } 
	   } 
	   } 
	 & & \\ 
 

\end{supertabular}
}

\end{document}
